\section*{Bab \ref{ch:b3:genetic-engineering} --- Rekayasa Genetika dan Bioteknologi}

\begin{solution}{exo:b3:genetic-engineering:1}
Sebuah titik asal replikasi, sehingga inangnya menyalin \omterm{def:b3:genetic-engineering:tools}{plasmidnya} lalu
meneruskannya kepada sel anaknya; sebuah \omterm{def:b3:genetic-engineering:tools}{penanda terpilih}, yang lazimnya
kekebalan antibiotik, sehingga hanya sel pembawa \omterm{def:b3:genetic-engineering:tools}{plasmidnya} yang tumbuh; dan
sebuah tapak restriksi tunggal (atau sekelompok tapak) tempat sisipannya
diligasikan tanpa memotong \omterm{def:b3:genetic-engineering:tools}{plasmidnya} di tempat lain.
\end{solution}

\begin{solution}{exo:b3:genetic-engineering:2}
Sebuah urutan berbasa enam tertentu muncul sekali tiap $4^{6} = 4096$ bp di
dalam DNA acak: sekitar $4.6\times 10^{6}/4096 \approx 1100$ fragmen yang
rata-rata \qty{4.1}{kb}. Genom yang sungguhan tidak acak --- komposisi basa,
pemakaian kodon, dan penghindaran sebagian urutan menggeser cacahnya (cacah
tapak EcoRI yang sesungguhnya pada \emph{E.~coli} beberapa ratus).
\end{solution}

\begin{solution}{exo:b3:genetic-engineering:3}
Denaturasi pada \qty{95}{\celsius} memisahkan untainya; penempelan pada
\qtyrange{50}{65}{\celsius} membiarkan \omterm{def:b3:genetic-engineering:pcr}{primernya} berpasangan dengan tapaknya;
pemanjangan pada \qty{72}{\celsius} membiarkan polimerasenya menyalin dari tiap
\omterm{def:b3:genetic-engineering:pcr}{primer}. Polimerase biasa akan musnah pada \qty{95}{\celsius}
pada daur pertamanya dan harus ditambahkan lagi tiga puluh kali; Taq
bertahan sepanjang seluruh daurnya.
\end{solution}

\begin{solution}{exo:b3:genetic-engineering:4}
Sebuah \omterm{def:b3:genetic-engineering:crispr}{RNA pemandu} yang kedua puluh basanya cocok dengan sasarannya, dan
sebuah \omterm{def:b3:genetic-engineering:crispr}{PAM} (NGG) tepat di sisi 3$'$ sasarannya pada untai bukan sasaran.
Sesudah \omterm{def:b3:dna-repair:lesions}{putus untai gandanya}: penyambungan ujung, yang meninggalkan sisipan
atau delesi kecil dan biasanya merusak kerangka baca gennya; atau \omterm{def:b3:dna-repair:dsb}{rekombinasi
homolog} dengan cetakan yang membawa urutan yang diinginkan, yang menuliskannya
ke dalamnya.
\end{solution}

\begin{solution}{exo:b3:genetic-engineering:5}
$50\times 1.9^{35} = 50\times 5.7\times 10^{9} \approx 3\times 10^{11}$
salinan. $\Delta C_{T} = 7.3$; nisbahnya $1.95^{7.3} = e^{7.3\ln 1.95}
\approx 130$: contoh pertamanya bercetakan kira-kira $130$ kali lebih banyak.
\end{solution}

\begin{solution}{exo:b3:genetic-engineering:6}
Bakteri tak dapat menyambung: salinan genom yang berintron akan
ditranskripsi menjadi pesan yang tak terpakai, sehingga cDNA yang bebas
intron, yang disalin dari mRNA matangnya, itulah yang dipakai. Rintangan lain:
promotor manusianya tidak dikenali polimerase bakteri (pasoklah promotor
bakteri dan tapak pengikat ribosom); pemakaian kodon manusia berbeda (sintesislah
gennya dengan kodon kegemaran inangnya); proteinnya mungkin tak melipat atau
menggumpal (turunkan suhunya, fusikan pada mitra yang larut, atau lipatlah
ulang dari badan inklusinya); glikosilasinya tak ada (pakailah khamir atau sel
mamalia bila hal itu penting); disulfidanya menuntut periplasma yang
mengoksidasi.
\end{solution}

\begin{solution}{exo:b3:genetic-engineering:7}
Sel punca yang tersasar disuntikkan ke dalam blastosista inang lalu bercampur
dengan sel inangnya sendiri, sehingga mencitnya terbangun dari dua genotipe ---
sebuah kimera --- dan hanya bila sebagian sel nutfahnya berasal dari sel yang
tersasar barulah ia meneruskan alel itu. Mengawinkan kimeranya dengan mencit
tipe liar memberi heterozigot (dari separuh gametnya bila separuh garis
nutfahnya tersasar); menyilangkan heterozigotnya satu sama lain memberi
seperempat homozigot tanpa gen itu.
\end{solution}

\begin{solution}{exo:b3:genetic-engineering:8}
Persis: $6.4\times 10^{9}\times 4^{-20}\times (1/16) \approx 4\times
10^{-4}$ tapak (GG milik \omterm{def:b3:genetic-engineering:crispr}{PAM-nya} memakan faktor $16$) --- tak satu pun, dalam
praktiknya. Dalam dua ketakcocokan terdapat $1 + 60 + 1710 = 1771$
urutan: $1771\times 4\times 10^{-4} \approx 0.7$ tapak kebetulan. Genom
yang sungguhan tidak acak, dan sebuah pemandu diperiksa langsung terhadapnya.
\end{solution}

\begin{solution}{exo:b3:genetic-engineering:9}
$p_{1} = 0.05 + 0.8\times 0.05\times 0.95 = 0.088$; $p_{2} = 0.152$;
$p_{3} = 0.255$; lalu $0.41$, $0.60$: kira-kira lima generasi untuk melewati
$0.5$, terhadap taksiran $\ln(1/(2\times 0.05))/\ln 1.8 \approx 4$.
\end{solution}

\begin{solution}{exo:b3:genetic-engineering:10}
Penyambungan ujung bekerja di setiap sel, termasuk sel punca yang mendiam,
yang di dalamnya \omterm{def:b3:dna-repair:dsb}{rekombinasi homolog}, sebuah proses S/G2, tak berdaya guna;
penyambungan itu tak menuntut cetakan untuk diantarkan dan hasilnya --- perusakan
penguatnya dengan cara apa pun --- sudah mengerjakan tugasnya, sedangkan
pembetulan menuntut urutan yang persis dan hanya memberi sebagian kecil
selnya. Lagi pula, suntingan yang sama berhasil bagi setiap mutasi
$\beta$-globin, sabit maupun talasemik. Kelemahannya: suntingan itu membuang
sebuah unsur pengatur (beserta peran lain apa pun yang dimilikinya),
meninggalkan alel sabitnya tetap di situ, dan campuran indelnya tak terkendali.
\end{solution}

\begin{solution}{exo:b3:genetic-engineering:11}
Dengan kebugaran homozigot $1 - s$: $p' = \bigl[p^{2}(1-s) + p(1-p)(1+c)
\bigr]/(1 - sp^{2})$. Dari kelangkaan, homozigotnya dapat diabaikan dan
penggeraknya tumbuh sebagai $(1+c)p$ berapa pun $s$; apakah ia mencapai
pemancangan diputuskan di dekat $p = 1$, tempat alel tipe liar yang langka
berfrekuensi $q$ kembali sebagai $q' \approx q(1-c)/(1-s)$: terpancang bila
$c > s$, kesetimbangan antara bila tidak. Kekebalannya: tiap konversi yang
gagal (penyambungan ujung alih-alih rekombinasi) meninggalkan alel yang tak
lagi dikenali pemandunya; bila alel semacam itu bugar --- sebuah indel
sekerangka di daerah yang tak esensial --- alel itu tak berongkos sedangkan
penggeraknya berongkos, sehingga seleksi memihaknya dan alel itu menggantikan
penggeraknya. Penawarnya adalah menyasar urutan esensial yang indelnya
mematikan, dan memakai beberapa pemandu sekaligus.
\end{solution}

\begin{solution}{exo:b3:genetic-engineering:12}
\omterm{def:b3:genetic-engineering:crispr}{Cas9} yang disuntikkan ke dalam sebuah zigot mungkin bekerja sesudah pembelahan
pertamanya, sehingga sel yang berbeda menerima perbaikan yang berbeda ---
mozaikisme; perbaikan tiap putusnya dipilih selnya sendiri, dan penyambungan
ujung memberi indel yang tak terduga serta delesi besar atau hilangnya
heterozigositas pada potongannya; hasil \omterm{def:b3:dna-repair:dsb}{rekombinasi homolognya} hanya sebagian
kecil; putus \omterm{prop:b3:genetic-engineering:editors}{luar sasaran} terjadi pada tapak yang tak semuanya dapat
diramalkan; dan embrionya hanya dapat diperiksa dengan mengurutkan beberapa
sel biopsi, yang tak dapat mewakili sisanya. Bukti yang dibutuhkan: metode
yang menyunting tiap selnya secara sama persis (sebelum fase S pertamanya)
dengan daya guna mendekati \qty{100}{\%}, pengurutan seluruh genom tiap sel
embrio hewan yang disunting yang memperlihatkan tak ada perubahan yang tak
disengaja, dan pemantauan jangka panjang hewan yang disunting sepanjang
beberapa generasi --- dan tak satu pun dari semua itu yang sudah ada.
\end{solution}

\begin{solution}{pb:b3:genetic-engineering:1}
\textbf{1.} Sekali tiap $4096$ bp. $P(\text{tanpa tapak}) = (1 - 1/4096)^{1500}
= e^{-0.37} = 0.69$.
\textbf{2.} Pilihlah enzim yang tak bertapak di dalam gennya; atau
perbanyaklah gennya dengan PCR memakai \omterm{def:b3:genetic-engineering:pcr}{primer} yang membawa tapak restriksi
baru pada ujung 5$'$-nya (atau pakailah pengklonan tumpul atau pengklonan tanpa ligasi).
\textbf{3.} \qty{0.1}{\micro g}\,$\times 2\times 10^{7} = 2\times
10^{6}$ koloni; \qty{10}{\%}-nya, $2\times 10^{5}$, membawa sisipannya.
\textbf{4.} Putih: \qty{10}{\%}. Separuh sisipannya berada pada orientasi
yang benar, sehingga $P(\text{tak satu pun dari } n \text{ koloni putih}) = 0.5^{n} \le 0.01$
memberi $n \ge 6.6$: pungutlah tujuh.
\textbf{5.} $32$ pelipatduaan, $2^{32} = 4.3\times 10^{9}$ sel, $2.1\times
10^{11}$ \omterm{def:b3:genetic-engineering:tools}{plasmid}; masing-masing $5500\times 650 = 3.6\times 10^{6}$ Da $=
5.9\times 10^{-18}$ g: \qty{1.3}{\micro g} \omterm{def:b3:genetic-engineering:tools}{plasmid}.
\textbf{6.} Pengklonan hanya menuntut replikasi, yang dipasok titik asal
\omterm{def:b3:genetic-engineering:tools}{plasmidnya}; promotor manusianya tak dibaca RNA polimerase bakteri, sehingga
bagi ekspresi sebuah promotor bakteri dan tapak pengikat ribosom harus
ditaruh di depan urutan penyandinya (yang bebas intron).
\textbf{7.} $N_{0} = N_{T}/2^{C_{T}}$: $10^{12}/2^{28} \approx 3700$
salinan; pada $C_{T} = 35$, $10^{12}/2^{35} \approx 29$ salinan.
\textbf{8.} $2^{n} = 10^{12}$: $n = 39.9$, empat puluh daur.
\textbf{9.} $10^{12}/2^{38} \approx 3.6$ salinan. Di atas kira-kira $40$
daur, satu molekul cemaran, atau artefak \omterm{def:b3:genetic-engineering:pcr}{primer}, mencapai
ambangnya, dan sebuah contoh yang kurang dari satu salinan tak bermakna.
\textbf{10.} Satu molekul melewatinya pada daur $40$: sebuah positif. Cegahlah
dengan ruang dan peralatan terpisah bagi penyiapan dan bagi penanganan
hasilnya, alur kerja satu arah, kendali negatif pada tiap jalan, serta
pemusnahan enzimatik amplikon yang terbawa.
\textbf{11.} $3700/0.4 \approx 9000$ salinan RNA.
\textbf{12.} Dilaporkan $2^{3.3} \approx 10$ kali lipat; sesungguhnya $1.9^{3.3} =
e^{3.3\ln 1.9} \approx 8.3$ kali lipat.
\textbf{13.} $40\times \qty{35}{\micro g} = \qty{1.4}{mg}$ tiap hari,
\qty{0.51}{g} tiap tahun; $\times 10^{7}$: \qty{5100}{kg} tiap tahun.
\textbf{14.} $5.1\times 10^{6}/5800 \approx 880$ mol; $5.3\times 10^{26}$
molekul.
\textbf{15.} $5.1\times 10^{6}$ g $/ (20$ g/L$) = 2.6\times 10^{5}$ L
$= \qty{260}{m^3}$ --- sepersepuluh kolamnya.
\textbf{16.} $5100\times 8000 = 4.1\times 10^{7}$ kg pankreas, pada
\qty{0.1}{kg} tiap ekor: $4\times 10^{8}$ ekor babi tiap tahun.
\textbf{17.} Hormon rekombinannya identik dengan hormon manusia, sehingga
tidak membangkitkan antibodi dan tidak menimbulkan reaksi alergi; hormon itu
tidak mengusung patogen hewan (virus, prion); pasokannya tidak bergantung pada
penyembelihan; dan urutannya dapat diperbaiki.
\textbf{18.} Residu yang menyentuh reseptornya tak berubah, sehingga
pengikatan dan pengisyaratannya adalah pengikatan dan pengisyaratan insulin;
residu yang diubah terletak di permukaan tempat molekul insulin berhimpun
menjadi dimer dan heksamer, dan mengubahnya mengubah kecepatan terurainya
depot suntikan itu menjadi monomer yang dapat diserap.
\textbf{19.} $6.4\times 10^{9}\times 4^{-20}/16 \approx 3.6\times
10^{-4}$: tak ada kecocokan kebetulan yang persis.
\textbf{20.} $1 + 60 + 1710 + 27\times 1140 = \num{32551}$ urutan;
$\num{32551}\times 3.6\times 10^{-4} \approx 12$ tapak \omterm{prop:b3:genetic-engineering:editors}{luar sasaran}
kebetulan.
\textbf{21.} $10^{7}\times 10^{-4} = 1000$ sel. Sebuah sel punca hidup dan
membelah seumur hidup pasiennya; sel yang telah kehilangan penekan
tumor, atau memperoleh sebuah translokasi, dapat mendirikan klon yang menjadi
leukemia.
\textbf{22.} $p_{1} = 0.0019$, $p_{2} = 0.0036$, $p_{3} = 0.0069$;
pertumbuhan geometrik sebesar $1.9$ tiap generasi memberi $\ln 500/\ln 1.9
\approx 10$; dengan mengiterasi rekursinya, frekuensinya melewati $0.5$ pada
generasi $11$ ($0.45 \to 0.68$).
\textbf{23.} Kira-kira setahun. Ongkos \qty{20}{\%} pada homozigotnya lebih
kecil daripada $c = 0.9$, sehingga penggeraknya tetap terpancang, lebih
lambat, sementara kebugaran rata-rata populasinya turun --- dan itulah tujuan
sebuah penggerak penekan; ongkos di atas $c$ akan menghentikannya pada
frekuensi antara.
\textbf{24.} $10^{6}\times 0.1 = 10^{5}$ alel kekebalan dalam satu
generasi. Karena tak dapat dipotong dan, bila bugar, tak berongkos sedangkan
penggeraknya berongkos, alel itu dipihak seleksi lalu menyebar, dan
penggeraknya tersingkir --- sehingga penggerak menyasar tapak yang hasil
penyambungan ujungnya mematikan, dengan beberapa pemandu sekaligus.
\textbf{25.} Sekitar $3700$ salinan di dalam contoh pertama; sekitar \qty{260}{m^3}
fermenter tiap tahun bagi insulin dunia; sekitar $11$ generasi, yakni setahun
nyamuk, bagi penggeraknya.
\end{solution}
